\section{Recursos do abntexto-ufc, elementos pós-textuais e revisão final}\label{sec:recursos-finais}

A classe reúne opções de perfil e módulos editoriais para que o mesmo projeto possa atender diferentes tipos de trabalho sem duplicar a lógica de formatação. Esta seção encerra o guia explicando os principais recursos, o fluxo de compilação/validação e os elementos que aparecem depois do texto principal.

\subsection{Perfis de documento}

A opção \texttt{type} seleciona o perfil acadêmico. A matriz de certificação atual exercita sete perfis: \texttt{undergraduate-capstone}, \texttt{specialization-capstone}, \texttt{masters-thesis}, \texttt{doctoral-thesis}, \texttt{research-project}, \texttt{anonymized-research-project} e \texttt{scientific-article}.

\guiapolitica{Perfis controlam a aplicabilidade de elementos e metadados. O objetivo é evitar que o usuário tenha de remover manualmente partes da classe para adaptar um TCC a uma dissertação, tese, projeto ou artigo.}

\guianormativa{ABNT NBR 15287:2025}{Os perfis \texttt{research-project} e \texttt{anonymized-research-project} usam a estrutura de projeto de pesquisa adotada pelo projeto. O perfil anonimizado adiciona uma política técnica de ocultação de dados pessoais para processos que exijam avaliação sem identificação; essa anonimização não é apresentada como regra geral da NBR 15287 \parencite{abnt15287}.}

\guiamecanismo{Selecione o perfil em \texttt{type}. Este TCC é o exemplo pedagógico completo do perfil \texttt{undergraduate-capstone}; o artigo possui fonte própria em \texttt{template/scientific-article.tex}. Projetos, dissertações, teses e especialização são exercitados pela matriz de perfis sem duplicar sete cópias deste guia.}

\guiavalidacao{Os gates \texttt{profiles}, \texttt{research-project} e as suítes específicas de artigo verificam a aplicabilidade e a composição dos perfis. A release também retém pares PDF/TeX dos sete perfis para inspeção visual.}

\subsection{Fontes e motores de compilação}

A opção \texttt{font=times|arial} seleciona a família textual desejada. A opção \texttt{strict-font} determina se a identidade literal é obrigatória ou se o ambiente pode usar fallback portátil.

\guiainstitucional{A política tipográfica UFC adotada pelo projeto trabalha com Arial ou Times New Roman. A certificação literal é executada em Windows com pdfLaTeX e LuaLaTeX; as fontes Microsoft não são redistribuídas pelo repositório.}

\guiapolitica{O documento de referência usa modo portátil para compilar em ambientes públicos. Uma compilação com fallback não deve ser confundida com evidência de identidade literal da família solicitada. A incorporação das fontes usadas no PDF é uma verificação separada, descrita na Seção~\ref{sec:formatacao-geral}.}

\guiamecanismo{A compilação padrão usa o motor configurado pelo fluxo do projeto. Para exercitar LuaLaTeX no repositório, a entrada de desenvolvimento é \texttt{make lua}; para a compilação padrão do documento canônico, use \texttt{make compile}.}

\subsection{Tabelas, códigos, algoritmos e equações}

O documento de referência exercita os módulos de objetos na Seção~\ref{sec:catalogo-exemplos}. Nem toda capacidade oferecida pela classe corresponde a uma exigência normativa específica.

\guiapolitica{Tabelas numéricas seguem a apresentação tabular exercitada pelo projeto e pelo gate IBGE. Já a escolha entre \texttt{listings} e \texttt{minted}, a apresentação de códigos e a numeração de linhas em código ou pseudocódigo são capacidades editoriais do projeto; não são declaradas como mandatos gerais da ABNT.}

\guiamecanismo{As chaves \texttt{tables}, \texttt{code} e \texttt{algorithms} selecionam módulos. O exemplo usa \texttt{tables=tabularray}, \texttt{code=listings} e \texttt{algorithms=algpseudocodex}. Códigos podem usar \texttt{ufclisting} ou \texttt{\string\ufcInputListing}; algoritmos usam \texttt{ufcalgorithm}; equações continuam usando os ambientes matemáticos do \LaTeX{}.}

\guiavalidacao{Os gates \texttt{table-ibge}, \texttt{code-typography}, \texttt{minted}, \texttt{algorithm-numbering}, \texttt{objects} e \texttt{math} protegem esses caminhos sem promover decisões editoriais de código/algoritmo a regras ABNT inexistentes.}

\subsection{PDF/A}

O cabeçalho de \texttt{main.tex} declara \texttt{pdfstandard=A-2b}. A suíte verifica o PDF gerado com veraPDF e também confere incorporação de fontes.

\guiainstitucional{O fluxo de depósito acompanhado pelo projeto exige PDF/A nos casos institucionais aplicáveis. As páginas atuais do Sistema de Bibliotecas devem ser verificadas conforme o tipo de trabalho \parencite{ufcdepositotcc2026,ufcdepositopos2026}.}

\guiapolitica{A escolha do perfil específico PDF/A-2b para o documento de referência e para a certificação é uma política técnica do projeto; ela não é apresentada como uma afirmação de que toda norma ABNT exige especificamente PDF/A-2b.}

\guiamecanismo{A declaração é feita por \texttt{\string\DocumentMetadata} antes de \texttt{\string\documentclass}. O usuário deve validar o PDF efetivamente gerado; a simples presença da declaração no fonte não prova conformidade do arquivo final.}

\guiavalidacao{A rota de release executa o gate PDF/A com veraPDF, além de validações de fontes e do validador do projeto.}

\subsection{Acessibilidade e limites de conformidade}

Acessibilidade em PDF envolve mais do que aparência visual. Metadados, estrutura, marcação, ordem de leitura e descrições alternativas podem exigir verificações automáticas e revisão humana dependendo do conteúdo e do nível de conformidade pretendido.

\guiapolitica{O \texttt{abntexto-ufc} valida somente as capacidades explicitamente cobertas por seus contratos. O projeto não declara conformidade PDF/UA completa apenas porque um PDF/A compila ou porque verificações estruturais selecionadas passam. Conteúdo alternativo adequado, leitura semântica de objetos e outros aspectos que não tenham prova automática permanecem sujeitos a revisão.}

\guiamecanismo{Mantenha metadados do documento coerentes e trate imagens/objetos com informação textual adequada à finalidade acadêmica. Não use um resultado verde do validador como substituto de revisão manual de conteúdo, ordem de leitura ou significado de descrições.}

\guiavalidacao{Os validadores classificam apenas checks que podem sustentar. Quando uma capacidade exige inspeção mais profunda ou humana, o resultado deve permanecer REVIEW/fora do PASS automático em vez de produzir uma falsa garantia.}

\subsection{Compilação reprodutível do projeto}

O \texttt{Makefile} é a entrada recomendada para evitar que o usuário precise memorizar a sequência de processadores auxiliares.

\guiamecanismo{Use \texttt{make compile} para o documento canônico. O alvo executa o motor LaTeX, chama Biber quando existe fonte bibliográfica ativa, chama \texttt{makeglossaries} quando há glossário, chama \texttt{makeindex} quando há índice e recompila o documento para estabilizar referências. \texttt{make check} executa o gate integrado de desenvolvimento; \texttt{make release-check} é reservado à matriz de release.}

\guiapolitica{Biber, glossário e índice são condicionais ao documento. Executar manualmente apenas uma passada de LaTeX e aceitar o primeiro PDF pode deixar referências, listas ou índice incompletos mesmo quando não há erro fatal.}

\subsection{Validador CLI/Deep}

O repositório fornece um validador local mais profundo em \texttt{tools/validate-ufc-pdf.py}. Ele trabalha sobre o PDF gerado e complementa a compilação com checks documentais disponíveis no ambiente CLI.

\guiamecanismo{Depois de compilar o trabalho, execute o validador sobre o PDF final. No repositório, o gate integrado usa essa ferramenta sobre \texttt{template/main.pdf}; para um documento do usuário, informe o caminho do PDF correspondente.}

\guiavalidacao{A suíte \texttt{pdf-validator} protege o fluxo CLI/Deep. Alguns checks dependem de ferramentas externas como Poppler e veraPDF; por isso a capacidade real do ambiente deve ser considerada ao interpretar o resultado.}

\subsection{Validador Web/Lite}

O diretório \texttt{validator/} contém a interface Web/Lite. Seu contrato é deliberadamente mais restrito que o CLI/Deep: o processamento do PDF ocorre localmente no navegador e checks que dependem de capacidades profundas não devem ser convertidos em PASS apenas porque a interface web não consegue executá-los.

\guiamecanismo{A interface Web/Lite é destinada à inspeção local de um PDF pelo navegador. Checks compartilhados usam os mesmos identificadores e semântica do contrato canônico quando a medição é equivalente; capacidades profundas permanecem identificadas como dependentes da rota CLI/Deep.}

\guiavalidacao{A release só considera essa superfície certificada quando um pacote estático limpo contém todas as dependências locais, o carregamento dos módulos passa e um PDF canônico real, além de um caso negativo, atravessa o caminho efetivo de análise do navegador. Testes sintéticos de schema/veredito, isoladamente, não substituem esse E2E.}

\subsection{Referências}

As referências aparecem imediatamente após o texto quando produzidas por \texttt{\string\ufcPrintReferences}. O banco bibliográfico é declarado com \texttt{\string\ufcAddBibliographyResource} e processado por Biber no fluxo recomendado.

\guianormativa{ABNT NBR 6023:2025}{Inclua na lista final as fontes de acordo com a política bibliográfica do trabalho e mantenha os dados necessários à identificação de cada documento \parencite{abnt6023}.}

\subsection{Glossário}

O glossário é um elemento opcional para definir termos especializados, siglas ou conceitos quando essa lista realmente auxiliar o leitor. O documento de referência o mantém habilitado para demonstrar a integração com o pacote \texttt{glossaries}.

\guiapolitica{Não use o glossário como duplicação automática de toda palavra técnica. Registre apenas termos cuja definição agregue valor ao trabalho e cuja repetição no corpo prejudicaria a leitura.}

\guiamecanismo{Selecione \texttt{glossary=glossaries}, declare as entradas no arquivo de glossário e produza o elemento com \texttt{\string\ufcPrintGlossary}. O alvo \texttt{make compile} chama o processador auxiliar quando necessário.}

\subsection{Apêndices}

Apêndices contêm material complementar elaborado pelo próprio autor e são identificados em sequência própria.

\guianormativa{estrutura pós-textual adotada pelo projeto}{Cada apêndice inicia em página própria. No escopo auditado, seu cabeçalho é centralizado, em letras maiúsculas, negrito e tamanho 12.}

\guiaexemplo{Este documento possui apêndices com material textual, questionário, código-fonte e orientação para documento externo autoral. Eles existem para exercitar diferentes comprimentos e conteúdos.}

\guiamecanismo{Declare cada apêndice com \texttt{\string\appendix\{Título\}} e mantenha o conteúdo em arquivo separado quando isso melhorar a organização do projeto.}

\subsection{Anexos}

Anexos são usados para material complementar de autoria externa ou não produzido pelo autor do trabalho.

\guianormativa{estrutura pós-textual adotada pelo projeto}{Cada anexo inicia em página própria e usa a mesma política de apresentação do cabeçalho exercitada para apêndices.}

\guiaexemplo{Se um documento externo em PDF precisar ser anexado, verifique legibilidade, direitos de uso, orientação de página e compatibilidade com o arquivo final. O exemplo deste projeto documenta a rota sem redistribuir material externo desnecessário.}

\guiamecanismo{Declare cada anexo com \texttt{\string\annex\{Título\}}. Arquivos externos continuam sujeitos a direitos de uso e às propriedades do PDF final; anexar um arquivo não transfere automaticamente sua licença para o projeto.}

\subsection{Índice}

O índice remissivo é diferente do sumário: ele organiza termos e assuntos para facilitar a localização de ocorrências no documento.

\guianormativa{ABNT NBR 6034:2004}{O índice é elemento opcional e possui norma específica de apresentação. A classe oferece o recurso com \texttt{imakeidx}, mas este TCC canônico deliberadamente não imprime o índice remissivo \parencite{abnt6034}.}

\guiamecanismo{Selecione \texttt{index=imakeidx}, marque termos com \texttt{\string\index} e produza o elemento com \texttt{\string\ufcPrintIndex}. O fluxo de compilação executa \texttt{makeindex} quando necessário.}

\subsection{Lombada}

\guianormativa{ABNT NBR 12225:2023}{A lombada possui aplicação condicional, normalmente relacionada à produção física do trabalho. O documento eletrônico de referência não a força em perfis nos quais não seja aplicável \parencite{abnt12225}.}

\guiapolitica{A ausência de uma lombada renderizada neste PDF eletrônico é deliberada. Quando houver produção física que exija lombada, a regra correspondente deve ser tratada no fluxo aplicável, sem transformar o arquivo eletrônico de referência em prova de uma característica física que ele não contém.}

\subsection{Checklist antes da entrega}

Antes de entregar um trabalho real produzido a partir deste modelo:

\begin{enumerate}
  \item revise todos os campos de \texttt{\string\ufcsetup} e confirme o perfil correto;
  \item mantenha apenas os pré-textuais aplicáveis e substitua conteúdos demonstrativos pelos dados reais;
  \item confirme citações e referências contra as fontes efetivamente consultadas;
  \item revise figuras, tabelas, quadros, códigos, algoritmos, equações, fontes e notas;
  \item verifique glossário, apêndices, anexos e índice somente quando pertinentes;
  \item compile pelo fluxo completo, incluindo processadores auxiliares condicionais;
  \item confirme tamanho A4, margens, paginação, tipografia e incorporação de fontes no PDF final;
  \item execute o validador local disponível e interprete corretamente checks automáticos, REVIEW e capacidades profundas;
  \item examine visualmente o PDF completo, inclusive quebras de página, listas, objetos e elementos externos;
  \item quando acessibilidade for requisito, revise também os aspectos semânticos/manuais que não possuem prova automática;
  \item confira as regras atuais do curso, programa e fluxo de depósito da UFC;
  \item para uma publicação do próprio template, aceite somente artefatos produzidos e certificados a partir do mesmo SHA de fonte.
\end{enumerate}

\guiapolitica{Um CI verde demonstra que o corpus testado respeitou os contratos executáveis do projeto. Ele não transforma automaticamente toda regra externa em prova formal de conformidade e não substitui a revisão acadêmica, institucional, visual ou de acessibilidade que permanecer manual.}

Este modelo comentado deve ser entendido como ponto de partida auditável: o usuário escreve o conteúdo acadêmico, enquanto a classe concentra a maior parte das decisões editoriais e o repositório registra a evidência usada para mantê-las.

\index{PDF/A}\index{acessibilidade}\index{validador}\index{glossário}\index{apêndice}\index{anexo}\index{índice}\index{projeto de pesquisa}\index{artigo científico}
